\chapter{الدوال الضربية الادعائية}
\label{ch:pretentious-multiplicative-functions}

\section*{نطاق الفصل وحالته}

يدرس هذا الفصل متوسطات الدوال الضربية المقيدة بالقيمة المطلقة بواحد من خلال قياس مدى تشبهها بالدوال الطورية \(n^{it}\)، ومع وجود بنية توافقية بالدوال 
\(\chi(n)n^{it}\). النواة الصارمة هي المسافة الادعائية وصيغة كمية من مبرهنة هالاش للمتوسطات الطويلة. أما الفترات القصيرة والارتباطات العليا ونتائج تشاو فتبقى خارج النواة.

\noindent\textbf{حالة الفصل:}
\textenglish{AUTHORED-DRAFT / NON-CITABLE}. فُتحت بوابة التأليف بعد مراجعة مستقلة ضيقة بحكم \textenglish{PASS-FOR-AUTHORING = YES}. تبقى جميع معرفات الفصل غير قابلة للاستشهاد حتى البناء والمراجعة المستقلة بعد التأليف واعتماد المالك.

\subsection*{حراس النطاق}
\begin{itemize}
  \item لا نفترض الضربية التامة إلا إذا صُرح بذلك.
  \item مبرهنة هالاش منقولة من المصادر وليست نتيجة مثبتة هنا.
  \item الصيغة الكمية المعتمدة تستعمل المجال \(|t|\le 2T\) وتحتفظ بالحد \(T^{-1/2}\).
  \item لا يُستنتج الإلغاء المتوسط من تباعد المقياس وحده مع تثبيت \(T\).
  \item نتائج الفترات القصيرة لا تُستخدم لإثبات أي نتيجة في النواة.
\end{itemize}

\section{الدوال الضربية والمتوسط الطويل}

\begin{definition}[الدالة الضربية والمتوسط المطبع]
\label{def:ch24-multiplicative-mean}
\resultid{ANT-DEF-24-01}
\provedhere
تسمى الدالة \(f:\mathbb N\to\mathbb C\) ضربـية إذا
\[
f(mn)=f(m)f(n)\qquad ((m,n)=1).
\]
ونفترض في هذا الفصل غالبًا أن \(|f(n)|\le1\). ويعرّف متوسطها الطويل المطبع بـ
\[
M_f(x)=\frac1x\sum_{n\le x}f(n).
\]
\end{definition}

السؤال المركزي هو: متى يكون \(M_f(x)\) صغيرًا؟ الصعوبة أن بعض الدوال الضربية لا تتذبذب عشوائيًا، بل قد تحاكي نموذجًا طوريًا منتظمًا. وللسياق الكلاسيكي لنظريات المتوسط ومعايير وجوده راجع عملَي Elliott \cite{elliott1975meanvalue,elliott1979probabilistic}.

\section{المسافة الادعائية}

\begin{definition}[المسافة الادعائية]
\label{def:ch24-pretentious-distance}
\resultid{ANT-DEF-24-02}
\citedresult[\cite{granvilleSoundararajan2008pretentious}]
لدالتين ضربيتين مقيدتين بالقيمة المطلقة بواحد نضع
\[
\mathbb D(f,g;x)^2
=
\sum_{p\le x}
\frac{1-\Re\!\left(f(p)\overline{g(p)}\right)}{p}.
\]
الجمع هنا على الأعداد الأولية فقط، والصيغة المربعة هي التطبيع الأساسي في هذا الفصل.
\end{definition}

إذا كانت قيم \(f(p)\) و\(g(p)\) متقاربة في المتوسط الموزون بـ\(1/p\)، كانت المسافة صغيرة. لكن الصفر في هذه المسافة حتى ارتفاع \(x\) لا يعني تطابق الدالتين على جميع الأعداد الطبيعية.

\begin{proposition}[خصائص المسافة الادعائية الأساسية]
\label{prop:ch24-distance-properties}
\resultid{ANT-PROP-24-01}
\provedhere
لدوال \(f,g,h\) مقيدة بواحد:
\[
\mathbb D(f,g;x)\ge0,
\qquad
\mathbb D(f,g;x)=\mathbb D(g,f;x),
\]
كما تحقق المتباينة المثلثية
\[
\mathbb D(f,h;x)
\le
\mathbb D(f,g;x)+\mathbb D(g,h;x).
\]
\end{proposition}

\begin{proof}
عدم السلبية والتماثل مباشران من التعريف. ولإثبات المتباينة المثلثية في قرص الوحدة كله، نستخدم رفعًا احتماليًا إلى دائرة الوحدة.

لكل \(z=re^{i\theta}\) حيث \(0\le r\le1\)، نعرّف متغيرًا عشوائيًا \(Z_z\) يأخذ القيمتين \(e^{i\theta}\) و\(-e^{i\theta}\) باحتمالين
\[
\frac{1+r}{2},
\qquad
\frac{1-r}{2},
\]
على الترتيب. عندئذ
\[
|Z_z|=1
\qquad\text{و}\qquad
\mathbb E Z_z=z.
\]

لكل أولي \(p\le x\)، نختار على فضاء احتمالي مشترك متغيرات مستقلة
\[
X_{f,p},\ X_{g,p},\ X_{h,p}
\]
من هذا النوع، بحيث
\[
\mathbb E X_{f,p}=f(p),
\quad
\mathbb E X_{g,p}=g(p),
\quad
\mathbb E X_{h,p}=h(p).
\]
وبالاستقلال وكون القيم أحادية القيمة المطلقة نحصل على
\[
\mathbb E\left|X_{f,p}-X_{g,p}\right|^2
=
2\left(1-\Re\!\left(f(p)\overline{g(p)}\right)\right).
\]
ومن ثم
\[
\mathbb D(f,g;x)^2
=
\sum_{p\le x}\frac{1}{2p}
\mathbb E\left|X_{f,p}-X_{g,p}\right|^2.
\]
أي أن \(\mathbb D(f,g;x)\) هو معيار الفرق بين المتجهين
\((X_{f,p})_{p\le x}\) و\((X_{g,p})_{p\le x}\)
في مجموع هيلبرت المباشر الموزون بالمعاملات \(1/(2p)\). لذلك تعطي متباينة مينكوفسكي
\[
\mathbb D(f,h;x)
\le
\mathbb D(f,g;x)+\mathbb D(g,h;x),
\]
وهو المطلوب.
\end{proof}

\section{النماذج الطورية ومقياس أفضل اقتراب}

الدالة \(n^{it}=e^{it\log n}\) ضربـية تمامًا وقيمتها المطلقة واحد. وهي النموذج الأساسي الذي قد تتظاهر به دالة ضربـية في المتوسط الطويل.

\begin{definition}[مقياس هالاش]
\label{def:ch24-halasz-measure}
\resultid{ANT-DEF-24-03}
\provedhere
لـ\(T\ge1\) نضع
\[
\mathcal M(f;x,T)
=
\min_{|t|\le2T}
\mathbb D(f,n^{it};x)^2.
\]
يقيس هذا العدد أفضل تشبه طوري للدالة \(f\) على الأوليات حتى \(x\).
\end{definition}

عند دراسة متتاليات حسابية أو التواءات توافقية تظهر النماذج \(\chi(n)n^{it}\)، لكن هذا توسيع إضافي لا يُخلط بالصورة غير التوافقية الأساسية.

\section{مبرهنة هالاش الكمية}

\begin{theorem}[هالاش: حد المتوسط الطويل]
\label{thm:ch24-halasz}
\resultid{ANT-THM-24-01}
\citedresult[\cite{halasz1968mittelwerte,granvilleSoundararajan2008pretentious,granvilleHarperSoundararajan2019halasz}]
لتكن \(f\) دالة ضربـية تحقق \(|f(n)|\le1\). عند \(x\ge2\) و\(T\ge1\)، لدينا بصيغة كمية معيارية
\[
\left|
\frac1x\sum_{n\le x}f(n)
\right|
\ll
(1+\mathcal M(f;x,T))e^{-\mathcal M(f;x,T)}
+
\frac1{\sqrt T},
\]
حيث الثابت المطلق لا يعتمد على \(f,x,T\).
\end{theorem}

المغزى أن المتوسط الكبير لا يستمر إلا إذا وجدت قيمة \(t\) في المجال \(|t|\le2T\) تجعل \(f\) قريبة من \(n^{it}\) على الأوليات. الحد الثاني يمنع تثبيت \(T\) ثم الادعاء بأن تباعد \(\mathcal M\) وحده يكفي.

\begin{corollary}[معيار الإلغاء المتوسط]
\label{cor:ch24-mean-cancellation}
\resultid{ANT-COR-24-01}
\provedhere
إذا أمكن اختيار دالة \(T=T(x)\to\infty\) بحيث
\[
\mathcal M(f;x,T(x))\to\infty,
\]
فإن
\[
\frac1x\sum_{n\le x}f(n)\to0.
\]
\end{corollary}

\begin{proof}
نطبق مبرهنة هالاش. يذهب الحد
\((1+\mathcal M)e^{-\mathcal M}\) إلى الصفر بتباعد \(\mathcal M\)، ويذهب \(T(x)^{-1/2}\) إلى الصفر لأن \(T(x)\to\infty\).
\end{proof}

\section{أمثلة أساسية}

\begin{remark}[أمثلة وتفسيرات]
\label{rem:ch24-examples}
\resultid{ANT-EX-24-01}
\begin{enumerate}
  \item \(f(n)=1\) تتظاهر تمامًا بالنموذج \(n^{i0}\)، ولذلك متوسطها يساوي واحدًا.
  \item \(f(n)=n^{it_0}\) تحقق مسافة صفر من النموذج نفسه متى كان \(|t_0|\le2T\).
  \item شخصية ديريشليه غير الرئيسية قد تحمل إلغاءً توافقيا، لكن تحليلها الكامل يستلزم فصل النموذج \(\chi(n)n^{it}\) عن النموذج الطوري وحده.
  \item دالتا موبيوس \(\mu\) وليوفيل \(\lambda\) تمثلان نموذجين مركزيين للتذبذب الضربي، لكن مبرهنة هالاش وحدها لا تمنح أفضل التقديرات المعروفة لمجاميعهما القصيرة.
\end{enumerate}
\end{remark}

\section{ماذا يعني الادعاء؟}

\begin{remark}[العائق الادعائي للمتوسط الكبير]
\label{rem:ch24-pretentious-obstruction}
\resultid{ANT-PRIN-24-01}
\citedresult[\cite{granvilleSoundararajan2008pretentious,granvilleHarperSoundararajan2019halasz}]
في نطاق المتوسطات الطويلة، العائق البنيوي الرئيس أمام الإلغاء هو التشبه بنموذج طوري مناسب. هذه صياغة تفسيرية منضبطة لعاقبة هالاش، وليست مبرهنة أوسع مستقلة عن شروطها.
\end{remark}

الفكرة لا تقول إن الدالة «تكذب»، بل إن سلوكها على الأوليات يقلد نموذجًا بسيطًا بما يكفي ليمنع الإلغاء المتوسط.

\section{حدود الاستدلال}

\begin{remark}[المتوسط الطويل ليس فترة قصيرة]
\label{rem:ch24-long-vs-short}
\resultid{ANT-PRIN-24-02}
\noindent\textenglish{METHODOLOGICAL-PRINCIPLE / INFERENCE-GUARDED}.
نجاح معيار هالاش على المجال \(n\le x\) لا يثبت إلغاءً منتظمًا في كل فترة \([x,x+h]\)، ولا يحدد تلقائيًا الارتباطات من نوع \(f(n)f(n+h)\). الانتقال إلى هذه المسائل يحتاج مدخلات مستقلة.
\end{remark}

هذا العنصر ليس مبرهنة جديدة ولا يحمل تصنيف \textenglish{PROVED-HERE}. وظيفته منع الانتقال غير المشروع من المتوسط العالمي إلى السلوك المحلي.

\section{الجبهة المفتوحة}

\begin{openproblem}[الفترات القصيرة والارتباطات]
\label{open:ch24-short-intervals-correlations}
\resultid{ANT-OPEN-24-01}
\openresult
تطوير قاموس الادعاء إلى متوسطات الفترات القصيرة والارتباطات العليا ومسائل تشاو يمثل جبهة مركزية. نتائج Matomäki--Radziwiłł \cite{matomakiRadziwill2016short} تبين أن المتوسطات القصيرة تحمل بنية عميقة إضافية، لكنها لا تدخل في برهان نواة هذا الفصل.
\end{openproblem}

\section*{خلاصة الفصل}

تضغط المسافة الادعائية سلوك الدالة الضربية على الأوليات في كمية هندسية موزونة. ويقيس \(\mathcal M(f;x,T)\) أقرب نموذج طوري، ثم تحول مبرهنة هالاش بعده عن جميع النماذج إلى إلغاء في المتوسط الطويل. غير أن الحد \(T^{-1/2}\) جزء حاكم من الصيغة، كما أن الانتقال إلى الفترات القصيرة والارتباطات يحتاج نظريات أخرى.

\section*{تمارين}
\begin{enumerate}
  \item احسب \(\mathbb D(1,n^{it};x)^2\) واكتبها كمجموع على الأوليات.
  \item تحقق أن \(\mathbb D(f,g;x)=0\) لا يفرض تطابق \(f\) و\(g\) بعد \(x\).
  \item استنتج معيار الإلغاء من مبرهنة هالاش مع توضيح ضرورة الشرطين على \(T(x)\) و\(\mathcal M\).
  \item قارن بين دور \(n^{it}\) ودور \(\chi(n)n^{it}\).
  \item اشرح لماذا لا تكفي نتيجة متوسط طويل لإثبات إلغاء في كل فترة قصيرة.
\end{enumerate}
